% ============================================================
\chapter{SQL --- Joins and Subqueries}
\label{ch:joins}
% ============================================================

The power of the relational model lies in its ability to link tables together. A student is enrolled in a department, which belongs to a university; each enrolment references a course, taught by a professor. To extract complete information---for example, ``which students take Professor Smith's courses?''---one must \emph{join} several tables. The join is the queen of SQL operations, and mastering it distinguishes the beginner from the expert. Subqueries, meanwhile, allow nesting questions within one another, offering considerable expressivity.

\section{Introduction}

\emph{Joins} combine data from multiple tables by exploiting the relationships
between them (foreign keys). \emph{Subqueries} (nested queries) allow using
the result of one query as input to another.

\section{Inner Join (INNER JOIN)}

\begin{definition}[Inner Join]
An \emph{inner join} combines rows from two tables where the join condition
is satisfied. Rows without a match are eliminated.
\end{definition}

\begin{codeblock}[title=\textbf{Equivalent Inner Join Syntaxes}]
\begin{minted}{sql}
-- Explicit syntax (recommended)
SELECT s.last_name, c.title, e.grade
FROM Students s
INNER JOIN Enrollments e ON s.student_id = e.student_id
INNER JOIN Courses c ON e.course_id = c.course_id;

-- Implicit syntax (old style, avoid)
SELECT s.last_name, c.title, e.grade
FROM Students s, Enrollments e, Courses c
WHERE s.student_id = e.student_id
  AND e.course_id = c.course_id;
\end{minted}
\end{codeblock}

\begin{codeoutput}
\begin{verbatim}
last_name | title      | grade
----------+------------+------
Brown     | Databases  | 88.5
Brown     | Algorithms | 76.0
Davis     | Databases  | 72.0
Davis     | Networks   | 91.0
Miller    | Algorithms | 95.5
Miller    | Statistics | 82.0
Wilson    | Databases  | 55.0
Taylor    | Statistics | 97.0
\end{verbatim}
\end{codeoutput}

\begin{bestpractice}[title=\textbf{Always Use Explicit JOIN Syntax}]
The implicit syntax (comma in \texttt{FROM}) is error-prone: forgetting a
condition in \texttt{WHERE} produces an accidental Cartesian product.
The \texttt{JOIN ... ON} syntax is clearer and separates join logic from filtering.
\end{bestpractice}

\section{Natural Join (NATURAL JOIN)}

\begin{codeblock}[title=\textbf{Natural Join}]
\begin{minted}{sql}
-- Joins on columns with the same name
SELECT * FROM Students NATURAL JOIN Enrollments;
\end{minted}
\end{codeblock}

\begin{warning}[title=\textbf{Danger of NATURAL JOIN}]
Natural join automatically determines common columns. If two tables share a
column name by coincidence (e.g., \texttt{last\_name} in both \texttt{Students}
and \texttt{Professors}), the result will be incorrect. Always prefer explicit
joins with \texttt{ON}.
\end{warning}

\section{Outer Joins (OUTER JOIN)}

\begin{definition}[Outer Join]
An \emph{outer join} preserves all rows from one table (or both), even those
without a match in the other table. Missing columns are filled with \texttt{NULL}.
\end{definition}

\subsection{LEFT JOIN}

\begin{codeblock}[title=\textbf{LEFT JOIN --- All Students, Even Without Enrollments}]
\begin{minted}{sql}
SELECT s.last_name, s.first_name, c.title, e.grade
FROM Students s
LEFT JOIN Enrollments e ON s.student_id = e.student_id
LEFT JOIN Courses c ON e.course_id = c.course_id
ORDER BY s.last_name;
\end{minted}
\end{codeblock}

If a student is not enrolled in any course, they still appear with
\texttt{NULL} for \texttt{title} and \texttt{grade}.

\subsection{RIGHT JOIN}

\begin{codeblock}[title=\textbf{RIGHT JOIN --- All Courses, Even Without Enrollments}]
\begin{minted}{sql}
SELECT c.title, s.last_name, e.grade
FROM Enrollments e
RIGHT JOIN Courses c ON e.course_id = c.course_id
LEFT JOIN Students s ON e.student_id = s.student_id;
\end{minted}
\end{codeblock}

\begin{remark}
\texttt{RIGHT JOIN} is less commonly used because it can always be rewritten
as a \texttt{LEFT JOIN} by reversing the table order.
SQLite does not support \texttt{RIGHT JOIN} (before version 3.39).
\end{remark}

\subsection{FULL OUTER JOIN}

\begin{codeblock}[title=\textbf{FULL OUTER JOIN}]
\begin{minted}{sql}
-- All students and all courses (PostgreSQL)
SELECT s.last_name, c.title
FROM Students s
FULL OUTER JOIN Enrollments e ON s.student_id = e.student_id
FULL OUTER JOIN Courses c ON e.course_id = c.course_id;
\end{minted}
\end{codeblock}

\begin{keyformulas}[title=\textbf{Join Types Summary}]
\begin{center}
\begin{tabular}{lp{8cm}}
\toprule
\textbf{Type} & \textbf{Description} \\
\midrule
\texttt{INNER JOIN} & Rows with matches in both tables \\
\texttt{LEFT JOIN} & All rows from the left table + matches \\
\texttt{RIGHT JOIN} & All rows from the right table + matches \\
\texttt{FULL OUTER JOIN} & All rows from both tables \\
\texttt{CROSS JOIN} & Cartesian product (all combinations) \\
\bottomrule
\end{tabular}
\end{center}
\end{keyformulas}

\section{Cross Join (CROSS JOIN)}

\begin{codeblock}[title=\textbf{Cartesian Product}]
\begin{minted}{sql}
-- All student-course combinations
SELECT s.last_name, c.title
FROM Students s
CROSS JOIN Courses c;
-- 5 students x 4 courses = 20 rows
\end{minted}
\end{codeblock}

\section{Self Join}

\begin{definition}[Self Join]
A \emph{self join} is a join of a table with itself. It requires aliases
to distinguish the two ``copies'' of the table.
\end{definition}

\begin{codeblock}[title=\textbf{Finding Pairs of Students in the Same Major}]
\begin{minted}{sql}
SELECT s1.last_name AS student1, s2.last_name AS student2, s1.major
FROM Students s1
JOIN Students s2 ON s1.major = s2.major
                 AND s1.student_id < s2.student_id;
\end{minted}
\end{codeblock}

\begin{codeoutput}
\begin{verbatim}
student1 | student2 | major
---------+----------+------
Brown    | Davis    | CS
Brown    | Wilson   | CS
Davis    | Wilson   | CS
Miller   | Taylor   | Math
\end{verbatim}
\end{codeoutput}

\section{Subqueries}

\begin{definition}[Subquery]
A \emph{subquery} is a \texttt{SELECT} statement nested inside another query.
It can appear in the \texttt{WHERE}, \texttt{FROM}, or \texttt{SELECT} clauses.
\end{definition}

\subsection{Scalar Subquery}

A subquery that returns exactly one value:

\begin{codeblock}[title=\textbf{Scalar Subquery}]
\begin{minted}{sql}
-- Students with a grade above the average
SELECT s.last_name, e.grade
FROM Students s
JOIN Enrollments e ON s.student_id = e.student_id
WHERE e.grade > (SELECT AVG(grade) FROM Enrollments);
\end{minted}
\end{codeblock}

\begin{codeoutput}
\begin{verbatim}
-- Average = 83.4375
last_name | grade
----------+------
Brown     | 88.5
Davis     | 91.0
Miller    | 95.5
Taylor    | 97.0
\end{verbatim}
\end{codeoutput}

\subsection{Subquery with IN}

\begin{codeblock}[title=\textbf{Subquery with IN}]
\begin{minted}{sql}
-- Students enrolled in a course taught by Prof. Smith
SELECT last_name, first_name
FROM Students
WHERE student_id IN (
    SELECT student_id
    FROM Enrollments
    WHERE course_id IN (
        SELECT course_id FROM Courses WHERE prof_id = 1
    )
);
\end{minted}
\end{codeblock}

\subsection{Correlated Subquery}

\begin{definition}[Correlated Subquery]
A subquery is \emph{correlated} when it references a column from the outer
query. It is re-evaluated for each row of the outer query.
\end{definition}

\begin{codeblock}[title=\textbf{Correlated Subquery with EXISTS}]
\begin{minted}{sql}
-- Students enrolled in at least one course
SELECT s.last_name, s.first_name
FROM Students s
WHERE EXISTS (
    SELECT 1
    FROM Enrollments e
    WHERE e.student_id = s.student_id
);

-- Students NOT enrolled in any course
SELECT s.last_name, s.first_name
FROM Students s
WHERE NOT EXISTS (
    SELECT 1
    FROM Enrollments e
    WHERE e.student_id = s.student_id
);
\end{minted}
\end{codeblock}

\begin{warning}[title=\textbf{Performance of Correlated Subqueries}]
Correlated subqueries are re-evaluated for each row. On large tables,
they can be slow. Modern optimizers often transform them into joins, but
it is good to know the alternative with \texttt{JOIN}:
\begin{minted}{sql}
-- Equivalent with LEFT JOIN + IS NULL
SELECT s.last_name FROM Students s
LEFT JOIN Enrollments e ON s.student_id = e.student_id
WHERE e.student_id IS NULL;
\end{minted}
\end{warning}

\subsection{Subquery in FROM (Derived Table)}

\begin{codeblock}[title=\textbf{Subquery in FROM}]
\begin{minted}{sql}
-- Average per student, then filter
SELECT sub.last_name, sub.avg_grade
FROM (
    SELECT s.last_name, AVG(e.grade) AS avg_grade
    FROM Students s
    JOIN Enrollments e ON s.student_id = e.student_id
    GROUP BY s.student_id, s.last_name
) AS sub
WHERE sub.avg_grade > 80;
\end{minted}
\end{codeblock}

\begin{codeoutput}
\begin{verbatim}
last_name | avg_grade
----------+----------
Brown     | 82.25
Davis     | 81.5
Miller    | 88.75
Taylor    | 97.0
\end{verbatim}
\end{codeoutput}

\subsection{Subquery in SELECT}

\begin{codeblock}[title=\textbf{Subquery in SELECT}]
\begin{minted}{sql}
SELECT s.last_name,
       (SELECT COUNT(*) FROM Enrollments e
        WHERE e.student_id = s.student_id) AS num_courses
FROM Students s;
\end{minted}
\end{codeblock}

\begin{codeoutput}
\begin{verbatim}
last_name | num_courses
----------+------------
Brown     | 2
Davis     | 2
Miller    | 2
Wilson    | 1
Taylor    | 1
\end{verbatim}
\end{codeoutput}

\section{ALL, ANY/SOME Operators}

\begin{codeblock}[title=\textbf{Comparison with ALL and ANY}]
\begin{minted}{sql}
-- Grade higher than ALL grades in course 101
SELECT s.last_name, e.grade
FROM Students s
JOIN Enrollments e ON s.student_id = e.student_id
WHERE e.grade > ALL (
    SELECT grade FROM Enrollments WHERE course_id = 101
);

-- Grade higher than at least ONE grade in course 101
SELECT s.last_name, e.grade
FROM Students s
JOIN Enrollments e ON s.student_id = e.student_id
WHERE e.grade > ANY (
    SELECT grade FROM Enrollments WHERE course_id = 101
);
\end{minted}
\end{codeblock}

\section{Common Table Expressions (CTE)}

\begin{definition}[CTE (Common Table Expression)]
A \emph{CTE} is a temporary named query, defined with the \texttt{WITH}
keyword, that can be referenced in the main query. It improves readability
and enables recursion.
\end{definition}

\begin{codeblock}[title=\textbf{Simple CTE}]
\begin{minted}{sql}
WITH Averages AS (
    SELECT student_id, AVG(grade) AS avg_grade
    FROM Enrollments
    GROUP BY student_id
)
SELECT s.last_name, a.avg_grade
FROM Students s
JOIN Averages a ON s.student_id = a.student_id
WHERE a.avg_grade > 80
ORDER BY a.avg_grade DESC;
\end{minted}
\end{codeblock}

\begin{codeblock}[title=\textbf{Recursive CTE --- Hierarchy}]
\begin{minted}{sql}
-- Categories table with parent reference
CREATE TABLE Categories (
    id INTEGER PRIMARY KEY,
    name TEXT,
    parent_id INTEGER REFERENCES Categories(id)
);

-- Recursive hierarchy traversal
WITH RECURSIVE Tree AS (
    -- Base case: roots
    SELECT id, name, parent_id, 0 AS depth
    FROM Categories WHERE parent_id IS NULL

    UNION ALL

    -- Recursive case
    SELECT c.id, c.name, c.parent_id, t.depth + 1
    FROM Categories c
    JOIN Tree t ON c.parent_id = t.id
)
SELECT * FROM Tree ORDER BY depth, name;
\end{minted}
\end{codeblock}

\section{Python Integration}

\begin{codeblock}[title=\textbf{Joins and Subqueries with Python}]
\begin{minted}{python}
import sqlite3

conn = sqlite3.connect('university.db')
cur = conn.cursor()

# Join: students and their courses
cur.execute('''
    SELECT s.last_name, c.title, e.grade
    FROM Students s
    JOIN Enrollments e ON s.student_id = e.student_id
    JOIN Courses c ON e.course_id = c.course_id
    WHERE e.grade > ?
    ORDER BY e.grade DESC
''', (80,))

print("Students with grade > 80:")
for name, course, grade in cur.fetchall():
    print(f"  {name:10} | {course:20} | {grade:.1f}")

# Subquery: students above average
cur.execute('''
    SELECT last_name, first_name FROM Students
    WHERE student_id IN (
        SELECT student_id FROM Enrollments
        WHERE grade > (SELECT AVG(grade) FROM Enrollments)
    )
''')
print("\nAbove average:", cur.fetchall())

conn.close()
\end{minted}
\end{codeblock}

\section*{Exercises}

\begin{exercise}
Write a query that displays each professor's name and the number of courses
they teach. Include professors who teach no courses.
\end{exercise}

\begin{exercise}
Find students enrolled in all courses taught by the professor with
\texttt{prof\_id} = 2. Use a subquery with \texttt{NOT EXISTS}.
\end{exercise}

\begin{exercise}
Write the same query as the previous exercise using \texttt{GROUP BY} and
\texttt{HAVING} with \texttt{COUNT}. Compare the two approaches.
\end{exercise}

\begin{exercise}
Find pairs of students who share at least two courses. Display the names
of both students and the number of shared courses.
\end{exercise}

\begin{exercise}
Using a recursive CTE, generate a table of numbers from 1 to 20 and
display their factorial.
\end{exercise}
